\section{Stationary torque residual and its nuisance offset}
\label{sec:stationary-residual}
Under a declared dq scaling and phase-system configuration, a reconstructed
flux map predicts
\begin{equation}
 M_{\mathrm{theo}}=\kt(\psid\iq-\psiq\id).
 \label{eq:theoretical-torque}
\end{equation}
Define measured minus theoretical torque,
\begin{equation}
 \boxed{\stationarytorqueresidual=M_{\mathrm{meas}}-M_{\mathrm{theo}}.}
 \label{eq:stationary-residual}
\end{equation}
The sign matters: a positive flux correction that increases theoretical torque
decreases this residual. The measured channel must be independent of the voltage
signals used for flux reconstruction; independence of electronics does not by
itself establish calibrated electromagnetic truth.

At stationary mechanical speed the schematic rotor balance is
\begin{equation}
 \rotorinertia\dot{\mechanicalspeed}
 =M_{\mathrm{em}}-\shafttorque-\mechanicallosstorque,
 \qquad \dot\omega_{\mathrm m}\simeq0.
 \label{eq:stationary-balance}
\end{equation}
Mechanical friction and windage, sensor-chain offset and some nearly constant
loss contribution can therefore separate measured shaft torque from the
electromagnetic model. Iron loss is not generally constant across a current
map even at fixed speed: flux amplitude, saturation and harmonics vary with
operating point. Current-dependent mechanical or electrical loss and calibration
gain can likewise create structure. A constant offset is consequently a first
model level, not a law of the machine.

For a group $g$ with approximately equal speed, thermal state, excitation and
measurement configuration, write
\begin{equation}
 \stationarytorqueresidual(i_d,i_q)
 =\torquegroupoffset+\delta r_M(i_d,i_q)+\varepsilon.
 \label{eq:residual-decomposition}
\end{equation}
The nuisance offset $c_g$ is not a target for flux correction. The structured
part $\delta r_M$ is the diagnostic signal; $\varepsilon$ represents measurement
and reconstruction scatter. The first hypothesis is
\begin{equation}
 \boxed{M_{\mathrm{meas}}-M_{\mathrm{theo,corr}}\simeq c_g
 \quad\text{within each qualified group}.}
 \label{eq:flat-residual-hypothesis}
\end{equation}
A nearly constant nonzero residual is therefore a successful first-order
consistency pattern. If spatially cross-validated low-order residual structure
remains, the next model level is a slowly varying nuisance surface. It should be
introduced only after the constant model fails, and it must not silently absorb
the same high-frequency structure that the flux diagnostic is meant to test.

Algebraically, centering removes one scalar per group. Geometrically, it aligns
the vertical level of the two torque surfaces without changing their shape.
Physically, it tolerates unmodeled mean loss and offset. Experimentally, it
prevents a large harmless bias from dominating a spatial-consistency metric.
